% !TEX program = xelatex
% Оформление: sadekovlar/misis-article-latex, см. LICENSE.
\documentclass{misisarticle}

\grnti{}
\asjc{}
\oecd{}
\title{Гиперсетевая генерация низкоранговых адаптеров для эффективной доменной адаптации VLA-моделей}
\misisauthor{А.~А.~Хабибуллин}{кафедра инженерной кибернетики НИТУ «МИСИС»}{Москва, Россия}{}

\setabstract{Рассматривается задача параметрически- и данно-эффективной доменной адаптации моделей vision-language-action (VLA). Систематизированы методы низкоранговой адаптации и генерации параметров гиперсетями. На основе метода Hyper-LoRA сформулирована схема адаптации компактной модели SmolVLA: гиперсеть порождает низкоранговые поправки к MLP-слоям языковой части замороженного визуально-языкового модуля, а эксперт действий обучается совместно с гиперсетью. Описаны математический аппарат, планируемые данные LIBERO и протокол проверки обобщения на LIBERO-Pro. Планируемый вклад связан с переносом Hyper-LoRA на SmolVLA, визуальным обусловливанием генератора и раздельной оценкой устойчивости к различным сдвигам распределения. Работа имеет аналитико-теоретический характер; экспериментальная проверка относится к следующему этапу исследования.}
\setkeywords{VLA, доменная адаптация, LoRA, гиперсети, SmolVLA, поведенческое клонирование, flow matching, LIBERO, LIBERO-Pro}
\hypersetup{pdftitle={Гиперсетевая генерация низкоранговых адаптеров для эффективной доменной адаптации VLA-моделей},pdfauthor={А. А. Хабибуллин}}

\begin{document}
\maketitle

\section{Введение}

Перенос идей фундаментальных моделей из обработки естественного языка и компьютерного зрения в робототехнику породил класс робототехнических фундаментальных моделей (RFM) и, в частности, VLA-моделей, которые по изображениям с камер и инструкции на естественном языке напрямую предсказывают низкоуровневые действия манипулятора~\cite{firoozi2023}. Такие модели обучаются поведенческим клонированием на демонстрациях; линия RT-1~\cite{rt1}, RT-2~\cite{rt2}, Open X-Embodiment~\cite{oxe}, Octo~\cite{octo} и OpenVLA~\cite{openvla} последовательно наращивала масштаб данных и общность, а компактная SmolVLA~\cite{smolvla} показала работоспособность подобной архитектуры при существенно меньших вычислительных затратах.

При переносе предобученной VLA в новый домен -- иную обстановку, набор задач или внешний вид объектов -- ее качество может падать из-за доменного разрыва, что требует доменной адаптации. Практически значима адаптация, одновременно параметрически-эффективная (мало обучаемых весов) и данно-эффективная (мало демонстраций). Остроту проблеме придает ограниченная обобщающая способность существующих политик: контролируемые возмущения бенчмарка LIBERO-Pro~\cite{liberopro} выявили существенное падение успешности моделей, показывающих высокие результаты на стандартном LIBERO~\cite{libero}. Это указывает на риск запоминания обучающих конфигураций и делает проверку обобщения ключевым требованием к методам адаптации.

Целью настоящей работы является теоретическая проработка метода гиперсетевой генерации низкоранговых адаптеров (Hyper-LoRA) для параметрически- и данно-эффективной доменной адаптации VLA-моделей, а также предварительная постановка задачи экспериментального исследования. На данном этапе исследование ограничено анализом предметной области, систематизацией методов и обоснованием архитектурных решений; практическая реализация и эмпирическая проверка отнесены к следующему этапу.

Для достижения поставленной цели решаются следующие задачи: проанализировать VLA-модели и проблему доменного разрыва; сопоставить низкоранговую адаптацию и гиперсетевую генерацию весов; описать механизм порождения адаптеров и согласование потоков как обучающий сигнал; обосновать выбор SmolVLA, LIBERO и LIBERO-Pro; сформулировать планируемый вклад и протокол проверки обобщения.

\section{Анализ работ}

\subsection{Робототехнические фундаментальные модели}

В рамках парадигмы VLA единая сеть принимает на вход изображения с камер робота и инструкцию на естественном языке, а на выходе напрямую предсказывает низкоуровневые действия манипулятора -- смещения и ориентацию схвата, состояние захвата и подобные управляющие величины. Восприятие, понимание языковой команды и порождение управляющего сигнала тем самым объединяются в одном дифференцируемом отображении.

Обучение VLA опирается на поведенческое клонирование: по набору экспертных демонстраций, каждая из которых представляет собой последовательность пар «наблюдение -- действие», модель настраивается воспроизводить действия эксперта при соответствующих наблюдениях и инструкции. Подобный подход не требует ни явной модели среды, ни функции вознаграждения, однако ставит качество в зависимость от объема и разнообразия обучающих траекторий.

Архитектурной основой развития VLA стал трансформер~\cite{attention}, механизм внимания которого позволяет обрабатывать разнородные модальности в общем представлении. Модель RT-1 продемонстрировала, что трансформерная политика, обученная на большом корпусе реальных демонстраций, способна исполнять множество инструкций одним манипулятором~\cite{rt1}. В RT-2 управление связано с предобученной визуально-языковой моделью (VLM), что перенесло семантические знания из веб-данных в робототехнические действия~\cite{rt2}. Инициатива Open X-Embodiment объединила демонстрации с разных платформ в общий корпус и обучила переносимые между воплощениями модели RT-X~\cite{oxe}.

Модель Octo предложила гибкую трансформерную политику с диффузионной головой генерации действий, допускающую разные конфигурации сенсоров и задач~\cite{octo}. OpenVLA продемонстрировала открытую VLA с практикой параметрически-эффективного дообучения под новые задачи~\cite{openvla}. Компактная модель SmolVLA объемом около 450 миллионов параметров развивает эту линию в сторону вычислительной доступности, сочетая VLM с отдельным экспертом действий и согласованием потоков (flow matching)~\cite{smolvla}. Малый размер, открытость и интеграция с экосистемой LeRobot делают ее удобной базой для исследования механизмов адаптации.

\subsection{Доменный разрыв и параметрически-эффективная адаптация}

Под доменом в настоящей работе понимается совокупность условий, в которых функционирует модель: обстановка и фон сцены, набор и внешний вид объектов, распределение начальных состояний, формулировки инструкций и состав целевых задач. Расхождение между обучающим и целевым распределениями составляет доменный разрыв; его порождают как сдвиги в визуальном облике сцены, так и изменения в семантике инструкций и в требуемом поведении.

Поведенческое клонирование усиливает чувствительность к доменному разрыву. Настроенная на воспроизведение конкретных демонстраций модель может опираться на статистические корреляции обучающего набора. Накопление ошибки вдоль траектории выводит исполнение в области пространства состояний, не покрытые демонстрациями, где предсказания особенно ненадежны.

Полная дообучка всех весов RFM требует хранения и оптимизации сотен миллионов параметров под каждый новый домен. При малом числе демонстраций это создает риск переобучения и катастрофического забывания. Параметрически-эффективная тонкая настройка (PEFT) сокращает такие издержки, замораживая основную модель и обучая лишь компактные дополнительные компоненты -- адаптеры, встраиваемые между слоями~\cite{houlsby}, либо низкоранговые поправки к существующим линейным преобразованиям~\cite{lora}.

Низкоранговая адаптация (LoRA) исходит из наблюдения, что необходимое для адаптации изменение весов часто допускает представление произведением двух матриц малой внутренней размерности~\cite{lora}. Базовый линейный слой при этом замораживается, а обучается лишь компактная низкоранговая поправка к нему. Для фиксированного адаптера поправку можно слить с исходной матрицей. При порождении адаптера по меняющемуся контексту стоимость генерации и применения необходимо учитывать отдельно.

Отдельная независимая настройка LoRA-адаптера под каждую задачу домена не использует в полной мере общую структуру близких задач. Естественное усиление состоит в том, чтобы разделять знание между задачами домена, порождая адаптеры из общего источника по описанию конкретной задачи. Альтернативный механизм обусловливания предлагает FiLM~\cite{film}: признаки промежуточных слоев модулируются аффинными преобразованиями, вычисляемыми по контексту. FiLM меняет активации, тогда как рассматриваемая схема порождает параметры низкоранговой поправки.

\subsection{Гиперсети как механизм генерации параметров}

В работе Ha и соавторов~\cite{hypernet} гиперсеть порождает веса другой, базовой сети, вместо того чтобы эти веса обучались напрямую. Гиперсеть отображает компактный вход -- например, эмбеддинг слоя или описание контекста -- в матрицы весов целевой сети. Тем самым параметризация базовой модели становится функцией входа гиперсети, а не фиксированным набором свободных переменных.

Когда на вход гиперсети подается контекст задачи -- инструкция на естественном языке, идентификатор условия или иные признаки текущего домена, -- порождаемые веса становятся специфичными для задачи. Контекстное обусловливание превращает гиперсеть в механизм, который вычисляет адаптер по описанию задачи, опираясь на усвоенную в ходе обучения общность между задачами. Энкодер контекста сжимает признаки задачи в эмбеддинг, а выходные головы преобразуют этот эмбеддинг в конкретные веса.

Анализ модульности гиперсетей~\cite{modularity} рассматривает выразительность и эффективность такого разложения. В мультизадачной настройке трансформеров общая гиперсеть может порождать параметры адаптеров сразу для множества задач~\cite{hyperformer}. Существенно, что общий генератор вынуждает разные задачи разделять структуру в пространстве весов. Возможность положительного переноса мотивирует применение этой схемы в робототехнике, но сама по себе не гарантирует обобщение на произвольные новые условия.

Сочетание двух идей дает компактную форму поправки к замороженной базе и способ порождать ее из контекста задачи. Вместо набора независимо обученных адаптеров обучается единая порождающая сеть, генерирующая низкоранговые веса для текущего контекста.

\subsection{Опорный метод Hyper-LoRA и наиболее близкие работы}

Опорным для настоящей работы выступает метод Hyper-LoRA~\cite{hyperlora}, в котором гиперсеть, обусловленная контекстом задачи, порождает специфичные для задачи параметры низкоранговой адаптации для робототехнической фундаментальной модели. Контекст задается инструкцией на естественном языке и кодируется энкодером в эмбеддинг фиксированной размерности, после чего выходные головы генератора порождают матрицы адаптера. Обучаемыми остаются параметры гиперсети, а не непосредственно порождаемые матрицы.

Эффективность метода обеспечивают два приема. Первый -- подача индекса слоя RFM как дополнительного входа гиперсети: единый набор выходных голов переиспользуется для разных слоев, а различие между слоями кодируется их индексом. Поскольку основная масса параметров гиперсети сосредоточена в выходных головах, такое переиспользование сокращает ее размер. Второй прием -- dropout контекстного эмбеддинга перед выходными головами, направленный на снижение переобучения генератора на конкретные обучающие задачи~\cite{hyperlora}.

Постановка в опорной работе строится на базе Octo с целевым доменом SIMPLER~\cite{simpler}: гиперсеть обучается на подмножестве задач, адаптеры добавляются к MLP-слоям, а диффузионная голова дополнительно дообучается. Обобщение проверяется без дополнительной настройки на оставшихся задачах домена~\cite{hyperlora}.

Для настоящей постановки наиболее релевантны Hyper-LoRA~\cite{hyperlora}, HyperVLA~\cite{hypervla} и SmolVLA~\cite{smolvla}. Первая задает опорный механизм генерации LoRA-весов, вторая исследует гиперсетевое порождение компактной политики с текстовым и визуальным контекстом, третья -- базовую архитектуру с экспертом действий и flow matching. В HyperVLA уже используются текстовые признаки T5 и визуальные признаки DINOv2, поэтому само включение визуальных признаков в контекст гиперсети не является новой идеей. Отличие настоящей постановки -- генерация именно LoRA-поправок для SmolVLA и сравнение вариантов контекста на LIBERO-Pro.

Настоящая постановка наследует базовый механизм Hyper-LoRA, но переносит его на SmolVLA. Визуально-языковой модуль замораживается, низкоранговые адаптеры порождаются поверх MLP-слоев его языковой части, а эксперт действий настраивается совместно с гиперсетью. Веса генерируются для каждого примера батча в процессе прямого прохода. Обобщение предполагается проверять на LIBERO-Pro без дообучения на его возмущенных условиях.

\subsection{Восприятие и перенос в смежных робототехнических задачах}

Смежные задачи восприятия для робототехники представлены в студенческом сборнике кафедры инженерной кибернетики НИТУ «МИСИС». В работе Д.~А.~Комаровой и А.~А.~Сыргулева~\cite{komarova} рассматривается определение дальности и детектирование объектов микроэлектроники в реальном времени. Работа Д.~Д.~Прохорова~\cite{prokhorov} посвящена определению геометрических размеров объектов на изображениях с использованием монокулярной оценки глубины. Эти исследования связаны с настоящей постановкой через визуальное восприятие сцены и использование нейросетевых представлений; генерация LoRA-адаптеров в них не рассматривается.

В работе RoboBenchMart~\cite{robobenchmart} непосредственно исследуется перенос VLA в среду розничного склада. Сравнение политик в новых конфигурациях сцен и с новыми объектами показывает значимость оценки за пределами условий обучения. Эта постановка близка настоящему исследованию проблемой доменного разрыва и необходимостью отдельно проверять обобщение после адаптации.

\section{Постановка задачи}

Под доменом понимается среда взаимодействия робота, заданная распределением визуальных сцен, набором манипулируемых объектов и пространством достижимых конфигураций манипулятора, вместе с конечным множеством задач, которые в этой среде требуется решать. Отдельная задача специфицируется инструкцией на естественном языке, описывающей желаемую цель манипуляции. Качество политики оценивается долей успешно завершенных эпизодов.

Поведение эксперта формализуется траекториями:
\begin{equation}
  \tau=(o_0,a_0,o_1,a_1,\ldots,o_H,a_H),
\end{equation}
где $o_t$ -- наблюдение в момент $t$, $a_t$ -- действие манипулятора, $H$ -- горизонт эпизода. Инструкция $c$ задается для всей траектории. Каждой задаче сопоставлен набор демонстраций, собранных под управлением эксперта.

Множество задач обозначим $\mathcal{T}$. Демонстрации доступны на обучающем подмножестве $\mathcal{T}_{\mathrm{train}}$, тогда как практический интерес представляет более широкое целевое множество условий $\mathcal{T}_{\mathrm{target}}$. Требуется по демонстрациям из $\mathcal{T}_{\mathrm{train}}$ построить параметрически- и данно-эффективную настройку предобученной VLA и оценить ее перенос на целевые условия, не представленные в обучении.

Пусть $\theta_0$ -- замороженные параметры VLM, $\phi$ -- параметры гиперсети, $\psi$ -- параметры эксперта действий. Требуется обучить $\phi$ и $\psi$, сохранив $\theta_0$ неизменными. Гиперсеть $H_\phi(c,l)$ порождает LoRA-поправку для слоя $l$; в расширенном варианте контекст $c$ дополнительно включает визуальные признаки текущей сцены.

Основная гипотеза состоит в том, что обусловленная контекстом генерация адаптеров позволит использовать общую структуру обучающих задач и улучшит перенос по сравнению с фиксированным адаптером при сопоставимых ресурсах. Проверке подлежат одновременно качество управления, потребность в демонстрациях, число обучаемых параметров и вычислительная стоимость. Преимущество метода на данном этапе не считается установленным.

\section{Описание методов}

\subsection{Поведенческое клонирование и согласование потоков}

Поведенческое клонирование формулирует обучение политики как задачу обучения с учителем: по наблюдению $o$ и инструкции $c$ предсказывается действие $a$, а целевыми метками служат действия эксперта из демонстраций. В общем виде минимизируется ожидаемое расхождение
\begin{equation}
  \mathcal{L}_{\mathrm{BC}}=\mathbb{E}_{(o,c,a)\sim\mathcal{D}}
  \bigl[d(\pi(o,c),a)\bigr].
\end{equation}
Здесь $\mathcal{D}$ -- распределение демонстрационных примеров, $\pi$ -- обучаемая политика, $d$ -- мера расхождения в пространстве действий. Такое обучение не требует взаимодействия со средой на этапе оптимизации.

Простейшая регрессия одного действия с квадратичной ошибкой плохо согласуется с многомодальностью демонстраций: цель может быть достижима несколькими допустимыми способами, а усреднение разных мод приводит к нежелательному «среднему» действию. Поэтому в современных VLA моделируется условное распределение действий, из которого порождаются выборки~\cite{octo,smolvla}.

В качестве генеративного механизма в SmolVLA применяется согласование потоков~\cite{smolvla,flowmatching}. Вводится семейство промежуточных состояний между шумовой выборкой $x_0$ и экспертным блоком действий $a$, параметризованное временем $s\in[0,1]$. Для прямолинейного пути
\begin{equation}
  x_s=(1-s)x_0+sa,\qquad u_s=a-x_0.
\end{equation}
Обучаемое векторное поле $v$ приближает $u_s$ при условии наблюдения и инструкции:
\begin{equation}
\begin{split}
  \mathcal{L}_{\mathrm{FM}}=\mathbb{E}_{o,c,a,x_0,s}\Bigl[
  \bigl\|v_{\theta_0,\phi,\psi}(x_s,s\mid o,c)\\
  {}-(a-x_0)\bigr\|_2^2\Bigr].
\end{split}
\label{eq:fm}
\end{equation}
Направление времени здесь выбрано от шума к данным; при реализации оно должно быть согласовано с соглашением используемого кода SmolVLA. На выводе интегрирование векторного поля преобразует шум в блок управляющих действий.

\subsection{Низкоранговая адаптация}

Для замороженного линейного слоя с весовой матрицей $W\in\mathbb{R}^{m\times n}$ исходное преобразование дополняется низкоранговой поправкой~\cite{lora}:
\begin{equation}
  h=Wx+\frac{\alpha}{r}W_{\mathrm{up}}W_{\mathrm{down}}x.
\label{eq:lora}
\end{equation}
Здесь $W_{\mathrm{down}}\in\mathbb{R}^{r\times n}$ проецирует вход в подпространство малой размерности $r$, а $W_{\mathrm{up}}\in\mathbb{R}^{m\times r}$ возвращает его в размерность выхода; $r\ll\min(m,n)$. Матрица $W$ остается неизменной. Множитель $\alpha/r$ задает масштаб поправки.

Полная настройка слоя потребовала бы переобучения $mn$ элементов, тогда как обычный LoRA-адаптер содержит $(m+n)r$ параметров. Отношение
\begin{equation}
  \frac{(m+n)r}{mn}
\end{equation}
мало при достаточно низком ранге. В гиперсетевой схеме эти матрицы порождаются, поэтому число оптимизируемых параметров всей системы определяется размером гиперсети и эксперта действий, а не только размером адаптера.

\subsection{Гиперсетевая параметризация адаптеров}

Гиперсеть устроена как композиция энкодера контекста и выходных голов. Сначала контекст $c$ отображается в скрытый эмбеддинг $e=f_\phi(c)$, затем головы $g_\phi$ порождают параметры адаптера. С учетом индекса слоя $l$ получаем
\begin{equation}
  \bigl(W_{\mathrm{down}}^{(l)},W_{\mathrm{up}}^{(l)}\bigr)
  =g_\phi\bigl(f_\phi(c),l\bigr).
\label{eq:hyper}
\end{equation}
При совпадающих размерностях адаптируемых матриц одни и те же выходные головы переиспользуются для разных слоев; для разных типов проекций необходимы совместимые по размерности головы. Такое разделение сокращает стоимость по сравнению с независимым полным генератором на каждый слой~\cite{hyperlora}.

Энкодер контекста может быть реализован многослойным перцептроном либо трансформером. В трансформерном варианте инструкция и представления индексов адаптируемых слоев подаются как токены единой последовательности, а внимание согласует контекст задачи с представлением конкретного слоя при порождении его весов.

В настоящей постановке адаптер порождается для каждого примера батча. На каждом прямом проходе по инструкции генерируются собственные матрицы $W_{\mathrm{down}}$ и $W_{\mathrm{up}}$, которые используются в MLP-слоях языковой части VLM. Градиент функции потерь~\eqref{eq:fm} проходит через замороженные слои и обновляет $\phi$ и $\psi$, не изменяя $\theta_0$.

Такое сквозное обучение проводится на фиксированных демонстрациях. Порождение весов внутри прямого прохода не предполагает обучение с подкреплением или сбор новых взаимодействий. Замораживание параметров VLM также не означает отключения вычисления градиента по входящим в него адаптерам.

\subsection{Визуальное обусловливание и регуляризация}

Способность гиперсети порождать осмысленные адаптеры для невиданных условий зависит от богатства контекста и риска переобучения. При бедном контексте выразительный генератор может запоминать соответствие «обучающая инструкция -- комплект весов». Для снижения этого риска предусматривается dropout контекстного эмбеддинга перед выходными головами~\cite{hyperlora}.

В расширенной конфигурации текстовый контекст предполагается дополнять визуальными признаками текущего наблюдения, извлекаемыми замороженным энкодером DINOv2~\cite{dinov2}. Тогда $c=(c_{\mathrm{text}},c_{\mathrm{visual}})$, и адаптер зависит не только от формулировки задачи, но и от наблюдаемой сцены. Сравнение текстового и текстово-визуального вариантов позволит проверить, помогает ли такая привязка при изменении объектов и их расположения.

Наличие вариаций в обучающем распределении не гарантирует перенос, а отсутствие прямых примеров определенного сдвига не доказывает его невозможность. Поэтому поведение по осям object, position, semantic и task рассматривается как предмет эмпирической проверки. Качество планируется представлять отдельно по каждой оси, без замены этих результатов единым средним показателем.

\section{Планируемый вклад в исследование}

Планируется следующий вклад:
\begin{enumerate}[leftmargin=0.65cm,itemsep=1pt,topsep=2pt]
  \item Реализация и исследование переноса Hyper-LoRA с Octo на компактную SmolVLA: генерация LoRA для MLP-слоев языковой части замороженного VLM и совместное обучение гиперсети с экспертом действий по flow matching.
  \item Исследование визуального обусловливания генератора с помощью DINOv2 и его сравнение с генерацией только по текстовой инструкции.
  \item Проверка параметрической и данной эффективности и устойчивости на LIBERO-Pro с раздельной оценкой по осям возмущений и контролем утечки тестовых условий.
\end{enumerate}

Идея гиперсетевой генерации LoRA не заявляется как новый метод сама по себе: она заимствована из~\cite{hyperlora}. Возможный вклад относится к архитектурной адаптации, способу задания контекста и воспроизводимому исследованию переноса. На текущем этапе сформулированы постановка и протокол; улучшение качества, экономия данных и вычислительная эффективность еще не подтверждены экспериментально.

\section{Наборы данных и протокол проверки}

\subsection{Обучающие демонстрации LIBERO}

В качестве источника данных выбран LIBERO~\cite{libero}. Полный бенчмарк включает 130 задач: LIBERO-Spatial, LIBERO-Object и LIBERO-Goal по 10 задач, а также LIBERO-100 со 100 задачами. Набор LIBERO-10 выделяет 10 наиболее длительных задач из LIBERO-100. В настоящей работе планируется стандартный протокол из четырех наборов Spatial, Object, Goal и LIBERO-10 -- всего 40 задач. Эти 40 задач следует отличать от полного объема LIBERO.

Для обучения планируется использовать готовую LeRobot-конверсию \texttt{physical-intelligence/libero}~\cite{liberodata}, примененную в работе SmolVLA. Карточка этой версии содержит 1693 эпизода и 273465 кадров на 40 задачах с частотой 10 кадров/с. Наблюдение включает RGB-изображения внешней и кистевой камер размером $256\times256\times3$, вектор состояния размерности 8 и инструкцию задачи. Действие имеет семь компонент: три компоненты поступательного перемещения, три компоненты вращения и команду захвата. Данные описывают симуляционную манипуляцию роботом Panda и представлены в формате LeRobot v2.0 с таблицами Parquet.

В исходном LIBERO опубликовано по 50 экспертных демонстраций на задачу, что для выбранных 40 задач дает 2000 траекторий до отбора. Это число относится к исходному набору HDF5 и не подменяет фактические 1693 эпизода выбранной LeRobot-конверсии. Различие версий необходимо учитывать при воспроизведении исследования.

Для исследования данной эффективности предлагается изменять бюджет демонстраций на задачу, например 5, 10, 25 и все доступные обучающие эпизоды, используя одинаковые подмножества для сравниваемых методов. Бюджеты уточняются с учетом реального количества эпизодов каждой задачи. Разбиение выполняется по целым траекториям, а не по отдельным кадрам. Подмножества для подбора гиперпараметров отделяются от обучающих демонстраций; тестовые эпизоды запускаются с независимыми начальными состояниями. При дополнительной проверке на удержанных задачах разбиение фиксируется на уровне задач.

На данном этапе наборы выбраны для последующей реализации: обучение, очистка и фактический отбор демонстраций еще не выполнены. Поэтому число реально использованных эпизодов и итоговые размеры разбиений будут зафиксированы после подготовки данных. Планируется использовать открытые веса SmolVLA и выполнить обучение на выбранной обучающей части LIBERO. При выборе уже дообученного checkpoint необходимо проверить состав его обучающих данных; идентификатор версии и состав предобучения должны быть сохранены для контроля утечки.

\subsection{LIBERO-Pro как проверка вне распределения}

LIBERO-Pro~\cite{liberopro} вводит контролируемые возмущения поверх LIBERO. В настоящем исследовании выбраны четыре группы: object -- изменения объектов; position -- изменения начальных положений; semantic -- перефразирование инструкций; task -- изменения целевой задачи. В актуальном репозитории бенчмарка дополнительно выделены изменения окружения (environment), которые в текущий протокол не включены. Выбранные тестовые условия предназначены для проверки переноса без дообучения на возмущенных сценах.

Демонстрации и результаты тестирования LIBERO-Pro не используются для обучения гиперсети, выбора checkpoint или подбора гиперпараметров. Для каждого варианта возмущений измеряется доля успешных эпизодов
\begin{equation}
  \mathrm{SR}=\frac{N_{\mathrm{success}}}{N_{\mathrm{episodes}}}.
\end{equation}
Результаты по каждой оси приводятся отдельно вместе с числом эпизодов, случайными seed и оценкой неопределенности. Отдельно фиксируются результаты на исходных условиях LIBERO.

\subsection{Сравниваемые конфигурации и воспроизводимость}

Намечено сравнение базовой SmolVLA без дополнительной адаптации, фиксированного LoRA-адаптера, гиперсети на тексте и гиперсети с визуальной привязкой. В конфигурациях с адаптацией режим обучения эксперта действий должен совпадать, чтобы его вклад не смешивался с эффектом генератора. Дополнительная конфигурация с обучением только эксперта действий позволит выделить этот эффект отдельно.

Помимо SR планируется учитывать число обучаемых параметров гиперсети и эксперта действий, размер порождаемого адаптера, объем демонстраций, время обучения, задержку генерации адаптера и память при выводе. Это необходимо, поскольку малый ранг адаптера сам по себе не гарантирует компактности всей гиперсетевой системы.

Средствами реализации выбраны PyTorch, LeRobot, SmolVLA, симуляционные среды LIBERO и LIBERO-Pro, а для визуальной привязки -- DINOv2. Версии кода, идентификаторы весов, разбиения, гиперпараметры и seed должны фиксироваться вместе с экспериментальными результатами.

\begin{thebibliography}{99}

\bibitem{firoozi2023}
Firoozi R., Tucker J., Tian S. et al. Foundation Models in Robotics: Applications, Challenges, and the Future. 2023. arXiv:2312.07843. \url{https://arxiv.org/abs/2312.07843}.

\bibitem{rt1}
Brohan A., Brown N., Carbajal J. et al. RT-1: Robotics Transformer for Real-World Control at Scale. 2022. arXiv:2212.06817. \url{https://arxiv.org/abs/2212.06817}.

\bibitem{rt2}
Brohan A., Brown N., Carbajal J. et al. RT-2: Vision-Language-Action Models Transfer Web Knowledge to Robotic Control. 2023. arXiv:2307.15818. \url{https://arxiv.org/abs/2307.15818}.

\bibitem{oxe}
Open X-Embodiment Collaboration. Open X-Embodiment: Robotic Learning Datasets and RT-X Models. 2023. arXiv:2310.08864. \url{https://arxiv.org/abs/2310.08864}.

\bibitem{octo}
Octo Model Team et al. Octo: An Open-Source Generalist Robot Policy. 2024. arXiv:2405.12213. \url{https://arxiv.org/abs/2405.12213}.

\bibitem{openvla}
Kim M. J., Pertsch K., Karamcheti S. et al. OpenVLA: An Open-Source Vision-Language-Action Model. 2024. arXiv:2406.09246. \url{https://arxiv.org/abs/2406.09246}.

\bibitem{smolvla}
Shukor M., Aubakirova D., Capuano F. et al. SmolVLA: A Vision-Language-Action Model for Affordable and Efficient Robotics. 2025. arXiv:2506.01844. \url{https://arxiv.org/abs/2506.01844}.

\bibitem{liberopro}
Zhou X., Xu Y., Tie G. et al. LIBERO-PRO: Towards Robust and Fair Evaluation of Vision-Language-Action Models Beyond Memorization. 2025. arXiv:2510.03827, v2, 2026. \url{https://arxiv.org/abs/2510.03827v2}.

\bibitem{libero}
Liu B., Zhu Y., Gao C. et al. LIBERO: Benchmarking Knowledge Transfer for Lifelong Robot Learning. \textit{Advances in Neural Information Processing Systems}. 2023. Vol. 36. P. 44776--44791. \url{https://arxiv.org/abs/2306.03310}.

\bibitem{attention}
Vaswani A., Shazeer N., Parmar N. et al. Attention Is All You Need. \textit{Advances in Neural Information Processing Systems}. 2017. P. 5998--6008. \url{https://arxiv.org/abs/1706.03762}.

\bibitem{houlsby}
Houlsby N., Giurgiu A., Jastrzebski S. et al. Parameter-Efficient Transfer Learning for NLP. \textit{Proceedings of ICML}. PMLR. 2019. Vol. 97. P. 2790--2799. \url{https://proceedings.mlr.press/v97/houlsby19a.html}.

\bibitem{lora}
Hu E. J., Shen Y., Wallis P. et al. LoRA: Low-Rank Adaptation of Large Language Models. \textit{ICLR}. 2022. \url{https://arxiv.org/abs/2106.09685}.

\bibitem{film}
Perez E., Strub F., de Vries H. et al. FiLM: Visual Reasoning with a General Conditioning Layer. \textit{Proceedings of AAAI}. 2018. Vol. 32, no. 1. \url{https://arxiv.org/abs/1709.07871}.

\bibitem{hypernet}
Ha D., Dai A., Le Q. V. HyperNetworks. \textit{ICLR}. 2017. \url{https://arxiv.org/abs/1609.09106}.

\bibitem{modularity}
Galanti T., Wolf L. On the Modularity of Hypernetworks. \textit{Advances in Neural Information Processing Systems}. 2020. Vol. 33. P. 10409--10419. \url{https://arxiv.org/abs/2002.10006}.

\bibitem{hyperformer}
Karimi Mahabadi R., Ruder S., Dehghani M., Henderson J. Parameter-Efficient Multi-Task Fine-Tuning for Transformers via Shared Hypernetworks. \textit{Proceedings of ACL-IJCNLP}. 2021. P. 565--576. \url{https://aclanthology.org/2021.acl-long.47/}.

\bibitem{hyperlora}
Sharma S., Xiong Z., Li K., Vuorio R., Whiteson S. Efficient Domain Adaptation of Robotic Foundation Models via Hypernetwork-Generated LoRA. \textit{NeurIPS 2024 Workshop}. 2024. \url{https://openreview.net/forum?id=SEvvWs3CAL}.

\bibitem{simpler}
Li X., Hsu K., Gu J. et al. Evaluating Real-World Robot Manipulation Policies in Simulation. 2024. arXiv:2405.05941. \url{https://arxiv.org/abs/2405.05941}.

\bibitem{hypervla}
Xiong Z., Li K., Wang Z., Jackson M., Foerster J., Whiteson S. HyperVLA: Efficient Inference in Vision-Language-Action Models via Hypernetworks. 2025. arXiv:2510.04898. \url{https://arxiv.org/abs/2510.04898}.

\bibitem{komarova}
Комарова Д. А., Сыргулев А. А. Создание системы определения дальности и детектирования объектов микроэлектроники в режиме реального времени. \textit{Искусственный интеллект в промышленных, коммерческих, медицинских и финансовых приложениях: сборник статей научно-технического семинара}. Вып. 5 / под ред. А. Р. Ефимова. М.: НИТУ «МИСИС», 2026. С. 77--81. \href{https://misis.ru/files/-/41f69c37b68def41ea566922c56b898e/kik_2025_2.pdf\#page=78}{Электронный сборник, статья с с.~77}.

\bibitem{prokhorov}
Прохоров Д. Д. Определение геометрических размеров объектов на изображениях с использованием монокулярной оценки глубины. \textit{Искусственный интеллект в промышленных, коммерческих, медицинских и финансовых приложениях: сборник статей научно-технического семинара}. Вып. 5 / под ред. А. Р. Ефимова. М.: НИТУ «МИСИС», 2026. С. 98--103. \href{https://misis.ru/files/-/41f69c37b68def41ea566922c56b898e/kik_2025_2.pdf\#page=99}{Электронный сборник, статья с с.~98}.

\bibitem{robobenchmart}
Soshin K., Krapukhin A., Spiridonov A., Bukhtuev G., Kuznetsov A., Shakhuro V., Shepelev D. RoboBenchMart: Benchmarking Robots in Retail Environment. 2025. arXiv:2511.10276, v2, 2026. \url{https://arxiv.org/abs/2511.10276v2}.

\bibitem{flowmatching}
Lipman Y., Chen R. T. Q., Ben-Hamu H. et al. Flow Matching for Generative Modeling. \textit{ICLR}. 2023. \url{https://arxiv.org/abs/2210.02747}.

\bibitem{dinov2}
Oquab M., Darcet T., Moutakanni T. et al. DINOv2: Learning Robust Visual Features without Supervision. 2023. arXiv:2304.07193. \url{https://arxiv.org/abs/2304.07193}.

\bibitem{liberodata}
Physical Intelligence. LIBERO: LeRobot dataset (physical-intelligence/libero). Карточка данных. \url{https://huggingface.co/datasets/physical-intelligence/libero}. Дата обращения: 07.11.2026.

\end{thebibliography}
\end{document}
